\documentclass[10pt,a4paper]{amsart}
\usepackage[margin=19mm]{geometry}
\usepackage{amsmath,amssymb,mathtools}
\usepackage[scr=rsfs,cal=euler]{mathalfa}
\usepackage{xfrac}
\usepackage[unicode,hidelinks,bookmarks=false]{hyperref}
\newcommand{\ie}{i.e.\ }
\newcommand{\cf}{cf.\ }
\newcommand{\resp}{respectively\ }
\newcommand{\Subsection}{\subsection}
\providecommand{\nobreakdash}{\nobreak\hbox{-}}
\setlength{\emergencystretch}{2em}
\multlinegap0pt
\usepackage{amssymb}
\newtheorem{theorem}{Theorem}
\title{Connectivities for $k$-knitted graphs and for minimal counterexamples to Hadwiger's Conjecture}
\author{Ken-ichi Kawarabayashi, Gexin Yu}
\date{Source: arXiv:2606.01586v2 (CC BY 4.0)}
\begin{document}
\maketitle

\noindent\emph{Source and licence.} Connectivities for $k$-knitted graphs and for minimal counterexamples to Hadwiger's Conjecture. Version arXiv:2606.01586v2 (CC BY 4.0), distributed by arXiv under the Creative Commons Attribution 4.0 International licence, http://creativecommons.org/licenses/by/4.0/, as recorded in the arXiv licence metadata for this exact version and shown on the abstract page https://arxiv.org/abs/2606.01586v2. Full licence notice: \texttt{licenses/arxiv-2606.01586-CC-BY-4.0.txt}.\par\smallskip

\noindent\textit{Editorial scope.} Counted unit: the article's theorem thm:unified (source0.tex lines 111--117). The article's complete original proof of it is delivered with it (source0.tex lines 119--138, one \texttt{proof} environment of 20 lines). No mathematics is written, repaired or re-derived: the statement and the proof are the article's own lines, in the article's own notation, and no retained line is substituted or re-flowed.  No further source-local context is needed: every cross-reference of every retained line resolves inside the two delivered spans.  Nothing was omitted from any retained span, so the package carries no omission record.  No retained line contains a citation, so the wrapper carries no bibliography and no bibliography entry.  No retained line contains a figure environment, an image-inclusion command or drawing code, so no image is delivered and the package declares no asset dependency.  Editorial formatting only, in addition: an AMS \texttt{amsart} preamble that restates the article's own declarations and macros used by the retained lines, namely the environments \texttt{theorem} on separate counters, together with the standard packages those lines need.  The retained environments print in this document as Theorem 1 in order of appearance, whereas the article prints the same items with the numbering of its own full document.  The complete native source of the article as distributed in the arXiv e-print for this exact version is bundled unchanged as \texttt{sources/201840/source0.tex}.
\begin{theorem}[General Reduction]\label{thm:unified}
Let $W$ be a graph, and let $S' \subseteq V(W)$ be a boundary set with $|S'| \ge 2$. Let $F$ be a forest in $W$ with $c \ge 1$ connected components such that $S' \subseteq V(F)$ and all leaves of $F$ belong to $S'$. Let $u \in V(W) \setminus V(F)$. Suppose that $d(u, F) \ge |S'| + k$, where $k = 2$ if $c = 1$, and $k = 1$ if $c \ge 2$. Then the induced subgraph $W[V(F) \cup \{u\}]$ contains a forest $F_0$ with $u \in V(F_0)$ such that $S' \subseteq V(F_0)$, all leaves of $F_0$ in $S'$, and at least one of the following holds:
\begin{enumerate}
    \item $|V(F_0)| < |V(F)|$ \quad \text{(Strict Order Reduction)}
    \item $|V(F_0)| = |V(F)|$ and $c(F_0) < c(F)$ \quad \text{(Strict Component Reduction)}
\end{enumerate}
\end{theorem}

\begin{proof}
Partition $S'$ across the components $T_1 \dots T_c$ of $F$ as $S'_i = S' \cap V(T_i)$. Since all leaves of $F$ are in $S'$, each $T_i$ is minimally spanned by $S'_i$. 
Let $N = N_W(u) \cap V(F)$. We have $|N| \ge |S'| + k$. Let $I$ be the index set of components containing at least one vertex of $N$. 

For each $i \in I$ and $s \in S'_i$, let $P_s$ be a shortest path in $T_i$ from $s$ to $N$. Let $y_s \in N$ be its terminus. Define $F_{\text{ess}}^{(i)} = \bigcup_{s \in S'_i} P_s$, and let $Y_i = \{y_s \mid s \in S'_i\}$. Let $Y = \bigcup Y_i$. Clearly, $|Y| \le |S'|$.

We construct an \textbf{intermediate graph} $F'_0$:
\[
F'_0 = \left( \bigcup_{i \in I} F_{\text{ess}}^{(i)} \right) \cup \{u\} \cup \{uy \mid y \in Y\} \cup \left( \bigcup_{j \notin I} T_j \right)
\]
While $F'_0$ connects the active components through $u$ and spans $S'$, it might contain cycles if paths in $F_{\text{ess}}^{(i)}$ coalesce and then diverge to multiple targets in $Y_i$. We resolve this by taking $F_0$ to be a \textbf{minimal spanning sub-forest} of $F'_0$ with respect to $S'$. By minimality, $F_0$ is cycle-free and all leaves of $F_0$ are in $S'$. 

We must guarantee $u \in V(F_0)$. If $u$ were removed during minimization, $F_0$ would be entirely contained within $F \setminus N$. This implies $\bigcup F_{\text{ess}}^{(i)}$ inherently spans $S'$, which forces $F_{\text{ess}}^{(i)} = T_i$ and $N = Y$. However, our threshold guarantees $|N| \ge |S'| + 1$, while $|Y| \le |S'|$, a contradiction. Thus, $u$ survives minimization.

Because $F_0 \subseteq F'_0$, any neighbor of $u$ in $N \setminus Y$ is completely excluded from $F_0$.

If $c = 1$, $k=2$. Thus $|N \setminus Y| \ge |S'|+2 - |S'| = 2$. The order strictly decreases: $|V(F_0)| \le |V(F)| - 2 + 1 < |V(F)|$, satisfying the requirement of our theorem. 

If $c \ge 2$, $k=1$. If $|I|=1$, then $d(u, T_i)\ge |S'|+1\ge |S_i'|+2$ for some sole component $T_i$. It follows that $W[T_i\cup \{u\}]$ has a smaller tree with desired properties. So $|I| \ge 2$.  Then the components merge into one component.  Thus $c(F_0) < c(F)$. Furthermore, $|V(F_0)| \le |V(F)| - |N \setminus Y| + 1 \le |V(F)|$.
\end{proof}

\end{document}
